\chapter{Isolated Singularities}
\label{ch:iv09}

An isolated singularity is classified by its Laurent principal part. Removable singularities have no negative powers, poles have finitely many, and essential singularities have infinitely many.

\section*{Learning goals}
After this chapter, the reader should be able to:
\begin{itemize}
\item classify isolated singularities as removable, poles, or essential using Laurent coefficients or growth criteria;
\item determine pole order from local factorizations;
\item apply removable-singularity theorems from boundedness or vanishing conditions;
\item use Casorati--Weierstrass or related behavior to recognize essential singularities;
\item construct local normal forms near poles and removable points;
\item distinguish singularity classification from global behavior of the function;
\end{itemize}
\section*{Conceptual roadmap}
\[
\boxed{\text{local holomorphic structure}\;\longrightarrow\;\text{integral or series representation}\;\longrightarrow\;\text{global consequence}.}
\]

\section{Removable singularities}
A bounded holomorphic function on a punctured neighborhood extends holomorphically across the puncture.

% BEGIN VOL04-EXPANSION IV09-example-01
\begin{example}[Removing a fake singularity]\label{ex:iv09-expand-01}
The quotient \(\sin z/z\) is undefined at zero in its displayed form, but
\[
\frac{\sin z}{z}=1-\frac{z^2}{3!}+\frac{z^4}{5!}-\cdots.
\]
The right side is analytic at zero and has value \(1\) there. Defining \(f(0)=1\) removes the singularity.
\end{example}
% END VOL04-EXPANSION IV09-example-01
\section{Poles}
A pole of order $m$ is characterized by $(z-z_0)^mf(z)$ extending holomorphically and nonzero at $z_0$.

\section{Essential singularities}
An essential singularity has infinitely many negative Laurent coefficients.

\section{Classification theorem}
Every isolated singularity is exactly one of removable, pole, or essential.

% BEGIN VOL04-EXPANSION IV09-example-02
\begin{example}[Detecting a pole and its order]\label{ex:iv09-expand-02}
For \(f(z)=e^z/z^3\), the numerator is nonzero at zero. Hence \(z^3f(z)=e^z\) extends holomorphically with nonzero value at zero. Therefore zero is a pole of order three. The leading Laurent term is \(z^{-3}\).
\end{example}
% END VOL04-EXPANSION IV09-example-02
\section{Riemann removable theorem}
Boundedness near the puncture is enough for removability.

\section{Casorati--Weierstrass}
Near an essential singularity the image is dense in the complex plane.

\section{Behavior of reciprocals}
Zeros of a holomorphic function correspond to poles of its reciprocal, with matching order.

% BEGIN VOL04-EXPANSION IV09-example-03
\begin{example}[An essential singularity]\label{ex:iv09-expand-03}
The Laurent series
\[
e^{1/z}=\sum_{n=0}^{\infty}\frac{1}{n!z^n}
\]
contains infinitely many negative powers. Thus zero is essential. No multiplication by a finite power of \(z\) can make the function bounded or holomorphic at the origin.
\end{example}
% END VOL04-EXPANSION IV09-example-03
\section{Singularity at infinity}
Use $w=1/z$ to classify behavior at infinity for functions on exterior domains.

\section{Core structural results}

\begin{theorem}[Classification by Laurent series]\label{thm:iv09-01}
An isolated singularity is removable iff all negative Laurent coefficients vanish, a pole iff only finitely many negative coefficients are nonzero, and essential otherwise.
\begin{proof}
The three cases correspond exactly to whether the Laurent series extends as a Taylor series, has a finite principal part with highest negative order, or has an infinite principal part.
\end{proof}
\end{theorem}

\begin{theorem}[Riemann removable singularity theorem]\label{thm:iv09-02}
If $f$ is holomorphic and bounded on a punctured disk around $z_0$, then $f$ extends holomorphically to $z_0$.
\begin{proof}
Use Cauchy's formula on a punctured annulus and estimate negative Laurent coefficients; boundedness forces every negative coefficient to vanish.
\end{proof}
\end{theorem}

\begin{theorem}[Pole-zero reciprocity]\label{thm:iv09-03}
$f$ has a zero of order $m$ at $z_0$ iff $1/f$ has a pole of order $m$ there.
\begin{proof}
Factor $f=(z-z_0)^mg$ with $g(z_0)\ne0$ and invert to obtain $(1/f)=(z-z_0)^{-m}(1/g)$ with holomorphic nonvanishing factor.
\end{proof}
\end{theorem}

\section{Worked examples}

\begin{example}[Removable singularities diagnostic]\label{ex:iv09-worked-01}
Give a rigorous worked analysis of Removable singularities. State the relevant holomorphic, contour, series, or singularity hypothesis and derive the central conclusion. A bounded holomorphic function on a punctured neighborhood extends holomorphically across the puncture. The decisive step is to use the complex-analytic structure globally rather than reason from real-variable intuition alone.
\end{example}

\begin{example}[Poles diagnostic]\label{ex:iv09-worked-02}
Give a rigorous worked analysis of Poles. State the relevant holomorphic, contour, series, or singularity hypothesis and derive the central conclusion. A pole of order $m$ is characterized by $(z-z_0)^mf(z)$ extending holomorphically and nonzero at $z_0$. The decisive step is to use the complex-analytic structure globally rather than reason from real-variable intuition alone.
\end{example}

\begin{example}[Essential singularities diagnostic]\label{ex:iv09-worked-03}
Give a rigorous worked analysis of Essential singularities. State the relevant holomorphic, contour, series, or singularity hypothesis and derive the central conclusion. An essential singularity has infinitely many negative Laurent coefficients. The decisive step is to use the complex-analytic structure globally rather than reason from real-variable intuition alone.
\end{example}

\begin{example}[Classification theorem diagnostic]\label{ex:iv09-worked-04}
Give a rigorous worked analysis of Classification theorem. State the relevant holomorphic, contour, series, or singularity hypothesis and derive the central conclusion. Every isolated singularity is exactly one of removable, pole, or essential. The decisive step is to use the complex-analytic structure globally rather than reason from real-variable intuition alone.
\end{example}

\section{Solved dossiers}
The dossiers are canonical solved problems. The provenance ledger separately records which retained corpus rules guided each topic and which dossiers were newly designed.

\begin{problem}[Removable singularities diagnostic]\label{prob:iv09-dossier-01}
Give a rigorous worked analysis of Removable singularities. State the relevant holomorphic, contour, series, or singularity hypothesis and derive the central conclusion.
\end{problem}
\begin{solution}
A bounded holomorphic function on a punctured neighborhood extends holomorphically across the puncture. The decisive step is to use the complex-analytic structure globally rather than reason from real-variable intuition alone.
\end{solution}

\begin{problem}[Poles diagnostic]\label{prob:iv09-dossier-02}
Give a rigorous worked analysis of Poles. State the relevant holomorphic, contour, series, or singularity hypothesis and derive the central conclusion.
\end{problem}
\begin{solution}
A pole of order $m$ is characterized by $(z-z_0)^mf(z)$ extending holomorphically and nonzero at $z_0$. The decisive step is to use the complex-analytic structure globally rather than reason from real-variable intuition alone.
\end{solution}

\begin{problem}[Essential singularities diagnostic]\label{prob:iv09-dossier-03}
Give a rigorous worked analysis of Essential singularities. State the relevant holomorphic, contour, series, or singularity hypothesis and derive the central conclusion.
\end{problem}
\begin{solution}
An essential singularity has infinitely many negative Laurent coefficients. The decisive step is to use the complex-analytic structure globally rather than reason from real-variable intuition alone.
\end{solution}

\begin{problem}[Classification theorem diagnostic]\label{prob:iv09-dossier-04}
Give a rigorous worked analysis of Classification theorem. State the relevant holomorphic, contour, series, or singularity hypothesis and derive the central conclusion.
\end{problem}
\begin{solution}
Every isolated singularity is exactly one of removable, pole, or essential. The decisive step is to use the complex-analytic structure globally rather than reason from real-variable intuition alone.
\end{solution}

\begin{problem}[Riemann removable theorem diagnostic]\label{prob:iv09-dossier-05}
Give a rigorous worked analysis of Riemann removable theorem. State the relevant holomorphic, contour, series, or singularity hypothesis and derive the central conclusion.
\end{problem}
\begin{solution}
Boundedness near the puncture is enough for removability. The decisive step is to use the complex-analytic structure globally rather than reason from real-variable intuition alone.
\end{solution}

\begin{problem}[Casorati--Weierstrass diagnostic]\label{prob:iv09-dossier-06}
Give a rigorous worked analysis of Casorati--Weierstrass. State the relevant holomorphic, contour, series, or singularity hypothesis and derive the central conclusion.
\end{problem}
\begin{solution}
Near an essential singularity the image is dense in the complex plane. The decisive step is to use the complex-analytic structure globally rather than reason from real-variable intuition alone.
\end{solution}

\begin{problem}[Behavior of reciprocals diagnostic]\label{prob:iv09-dossier-07}
Give a rigorous worked analysis of Behavior of reciprocals. State the relevant holomorphic, contour, series, or singularity hypothesis and derive the central conclusion.
\end{problem}
\begin{solution}
Zeros of a holomorphic function correspond to poles of its reciprocal, with matching order. The decisive step is to use the complex-analytic structure globally rather than reason from real-variable intuition alone.
\end{solution}

\begin{problem}[Singularity at infinity diagnostic]\label{prob:iv09-dossier-08}
Give a rigorous worked analysis of Singularity at infinity. State the relevant holomorphic, contour, series, or singularity hypothesis and derive the central conclusion.
\end{problem}
\begin{solution}
Use $w=1/z$ to classify behavior at infinity for functions on exterior domains. The decisive step is to use the complex-analytic structure globally rather than reason from real-variable intuition alone.
\end{solution}

\begin{problem}[Classification by Laurent series proof dossier]\label{prob:iv09-dossier-09}
Prove the following theorem-level statement and identify the essential hypothesis: An isolated singularity is removable iff all negative Laurent coefficients vanish, a pole iff only finitely many negative coefficients are nonzero, and essential otherwise.
\end{problem}
\begin{solution}
The three cases correspond exactly to whether the Laurent series extends as a Taylor series, has a finite principal part with highest negative order, or has an infinite principal part. This identifies the mechanism that makes the complex-analytic conclusion stronger than its real-variable analogue.
\end{solution}

\begin{problem}[Riemann removable singularity theorem proof dossier]\label{prob:iv09-dossier-10}
Prove the following theorem-level statement and identify the essential hypothesis: If $f$ is holomorphic and bounded on a punctured disk around $z_0$, then $f$ extends holomorphically to $z_0$.
\end{problem}
\begin{solution}
Use Cauchy's formula on a punctured annulus and estimate negative Laurent coefficients; boundedness forces every negative coefficient to vanish. This identifies the mechanism that makes the complex-analytic conclusion stronger than its real-variable analogue.
\end{solution}

\begin{problem}[Pole-zero reciprocity proof dossier]\label{prob:iv09-dossier-11}
Prove the following theorem-level statement and identify the essential hypothesis: $f$ has a zero of order $m$ at $z_0$ iff $1/f$ has a pole of order $m$ there.
\end{problem}
\begin{solution}
Factor $f=(z-z_0)^mg$ with $g(z_0)\ne0$ and invert to obtain $(1/f)=(z-z_0)^{-m}(1/g)$ with holomorphic nonvanishing factor. This identifies the mechanism that makes the complex-analytic conclusion stronger than its real-variable analogue.
\end{solution}

\begin{problem}[Chapter synthesis]\label{prob:iv09-dossier-12}
Connect the definitions and principal theorems of this chapter into one reusable complex-analysis strategy.
\end{problem}
\begin{solution}
An isolated singularity is classified by its Laurent principal part. Removable singularities have no negative powers, poles have finitely many, and essential singularities have infinitely many. The reusable strategy is to identify the analytic domain, choose the correct local expansion or contour representation, apply the structural theorem, and then audit singularities and topology.
\end{solution}

\section{Exercises with complete solutions}

\begin{exercise}\label{exr:iv09-01}
State and apply the central principle behind Removable singularities. Give one concrete consequence.
\end{exercise}
\begin{hint}
If all negative Laurent coefficients vanish, the singularity is removable; if finitely many occur, it is a pole; infinitely many signal essential behavior.
\end{hint}
\begin{solution}
A bounded holomorphic function on a punctured neighborhood extends holomorphically across the puncture. The same principle can then be reused in later contour, residue, conformal, or Riemann-surface arguments.
\end{solution}

\begin{exercise}\label{exr:iv09-02}
State and apply the central principle behind Poles. Give one concrete consequence.
\end{exercise}
\begin{hint}
For a pole, factor \(f(z)=(z-a)^{-m}g(z)\) with \(g(a)\ne0\) and identify the smallest valid \(m\).
\end{hint}
\begin{solution}
A pole of order $m$ is characterized by $(z-z_0)^mf(z)$ extending holomorphically and nonzero at $z_0$. The same principle can then be reused in later contour, residue, conformal, or Riemann-surface arguments.
\end{solution}

\begin{exercise}\label{exr:iv09-03}
State and apply the central principle behind Essential singularities. Give one concrete consequence.
\end{exercise}
\begin{hint}
A bounded punctured-neighborhood function has a removable singularity; extend it by the Cauchy-limit value.
\end{hint}
\begin{solution}
An essential singularity has infinitely many negative Laurent coefficients. The same principle can then be reused in later contour, residue, conformal, or Riemann-surface arguments.
\end{solution}

\begin{exercise}\label{exr:iv09-04}
State and apply the central principle behind Classification theorem. Give one concrete consequence.
\end{exercise}
\begin{hint}
Use \((z-a)^mf(z)\to c\ne0\) to certify a pole of order \(m\).
\end{hint}
\begin{solution}
Every isolated singularity is exactly one of removable, pole, or essential. The same principle can then be reused in later contour, residue, conformal, or Riemann-surface arguments.
\end{solution}

\begin{exercise}\label{exr:iv09-05}
State and apply the central principle behind Riemann removable theorem. Give one concrete consequence.
\end{exercise}
\begin{hint}
For essential singularities, inspect the Laurent principal part or invoke Casorati--Weierstrass only after ruling out removable and pole cases.
\end{hint}
\begin{solution}
Boundedness near the puncture is enough for removability. The same principle can then be reused in later contour, residue, conformal, or Riemann-surface arguments.
\end{solution}

\begin{exercise}\label{exr:iv09-06}
State and apply the central principle behind Casorati--Weierstrass. Give one concrete consequence.
\end{exercise}
\begin{hint}
Classification is local: ignore remote poles and zeros when analyzing the singularity at \(a\).
\end{hint}
\begin{solution}
Near an essential singularity the image is dense in the complex plane. The same principle can then be reused in later contour, residue, conformal, or Riemann-surface arguments.
\end{solution}

\begin{exercise}\label{exr:iv09-07}
State and apply the central principle behind Behavior of reciprocals. Give one concrete consequence.
\end{exercise}
\begin{hint}
Zeros of the reciprocal correspond to poles of the original function; use this to switch perspectives when convenient.
\end{hint}
\begin{solution}
Zeros of a holomorphic function correspond to poles of its reciprocal, with matching order. The same principle can then be reused in later contour, residue, conformal, or Riemann-surface arguments.
\end{solution}

\begin{exercise}\label{exr:iv09-08}
State and apply the central principle behind Singularity at infinity. Give one concrete consequence.
\end{exercise}
\begin{hint}
Check whether cancellation removes an apparent singularity before classifying the unreduced formula.
\end{hint}
\begin{solution}
Use $w=1/z$ to classify behavior at infinity for functions on exterior domains. The same principle can then be reused in later contour, residue, conformal, or Riemann-surface arguments.
\end{solution}

\section*{Chapter summary}
The chapter now forms part of the canonical complex-analysis chain from holomorphicity through residues, global continuation, and Riemann surfaces.
% BEGIN VOL04-EXPANSION IV09-exercises-01
\section{Graded supplementary exercises}
These exercises extend the protected chapter with a balanced set of computations, proofs, hypothesis tests, applications, and challenges.

\subsection*{Standard computations}

\begin{exercise}[Removable quotient]\label{exr:iv09-expand-01}
Classify the singularity of \((e^z-1)/z\) at zero.
\end{exercise}
\begin{hint}
Expand the numerator or take a limit.
\end{hint}
\begin{solution}
The quotient tends to one and has a Taylor expansion beginning \(1+z/2+\cdots\), so the singularity is removable.
\end{solution}

\begin{exercise}[Pole order]\label{exr:iv09-expand-02}
Classify zero for \((1+z)/z^5\).
\end{exercise}
\begin{hint}
The numerator is nonzero at zero.
\end{hint}
\begin{solution}
Zero is a pole of order five.
\end{solution}

\begin{exercise}[Essential exponential]\label{exr:iv09-expand-03}
Classify zero for \(e^{1/z^2}\).
\end{exercise}
\begin{hint}
Expand the exponential.
\end{hint}
\begin{solution}
The Laurent series has infinitely many negative powers, so zero is essential.
\end{solution}

\begin{exercise}[Zero versus pole]\label{exr:iv09-expand-04}
Classify zero for \(z^3/(e^z-1)^2\).
\end{exercise}
\begin{hint}
Use \(e^z-1=z+O(z^2)\).
\end{hint}
\begin{solution}
The denominator has order two and the numerator order three, leaving a zero of order one after cancellation; the apparent singularity is removable.
\end{solution}

\begin{exercise}[Limit criterion]\label{exr:iv09-expand-05}
If \(\lim_{z\to a}(z-a)^2f(z)=7\ne0\), what singularity does \(f\) have at \(a\)?
\end{exercise}
\begin{hint}
Use the pole characterization.
\end{hint}
\begin{solution}
The point is a pole of order two.
\end{solution}

\subsection*{Proofs}

\begin{exercise}[Bounded removable theorem]\label{exr:iv09-expand-06}
Prove that a bounded holomorphic function on a punctured disk has a removable singularity at the center.
\end{exercise}
\begin{hint}
Apply Cauchy's formula on circles avoiding the point or examine Laurent coefficients.
\end{hint}
\begin{solution}
For negative Laurent coefficients, the coefficient estimate contains a positive power of the circle radius. Boundedness forces each such coefficient to zero as the radius shrinks, leaving a Taylor series.
\end{solution}

\begin{exercise}[Pole characterization]\label{exr:iv09-expand-07}
Prove that \(a\) is a pole of order \(m\) exactly when \((z-a)^m f(z)\) extends holomorphically and nonvanishingly at \(a\).
\end{exercise}
\begin{hint}
Use the finite principal part factorization.
\end{hint}
\begin{solution}
A pole of order \(m\) has Laurent form \((z-a)^{-m}g(z)\) with \(g(a)\ne0\). Multiplying removes exactly that pole; the converse follows by dividing the nonvanishing extension by \((z-a)^m\).
\end{solution}

\begin{exercise}[Removable from finite limit]\label{exr:iv09-expand-08}
Show that a finite limit of \(f(z)\) as \(z\to a\) makes the isolated singularity removable.
\end{exercise}
\begin{hint}
A finite limit gives local boundedness.
\end{hint}
\begin{solution}
Near \(a\), the function is bounded, so the removable singularity theorem applies. Defining the missing value to be the limit yields continuity and holomorphicity.
\end{solution}

\begin{exercise}[Essential criterion]\label{exr:iv09-expand-09}
Show that an isolated singularity is essential exactly when its Laurent principal part has infinitely many nonzero terms.
\end{exercise}
\begin{hint}
Use the exhaustive Laurent classification.
\end{hint}
\begin{solution}
No negative terms gives a removable point; finitely many gives a pole. The only remaining case is infinitely many negative terms, which therefore characterizes an essential singularity.
\end{solution}

\subsection*{Counterexamples and hypothesis tests}

\begin{exercise}[Unbounded means pole]\label{exr:iv09-expand-10}
Does unboundedness near an isolated singularity always imply a pole?
\end{exercise}
\begin{hint}
Recall essential singularities.
\end{hint}
\begin{solution}
No. Essential singularities are also unbounded in every punctured neighborhood. One needs finite principal part or equivalent pole behavior.
\end{solution}

\begin{exercise}[Vanishing product]\label{exr:iv09-expand-11}
If \((z-a)f(z)\to0\), must \(f\) have a removable singularity?
\end{exercise}
\begin{hint}
Apply the removable theorem first to \(g(z)=(z-a)f(z)\).
\end{hint}
\begin{solution}
Yes. The function \(g\) extends holomorphically with \(g(a)=0\), so \(g(z)=(z-a)h(z)\) for a holomorphic \(h\). Hence \(f=h\) on the punctured disk and extends holomorphically across \(a\).
\end{solution}

\begin{exercise}[Classification requires isolation]\label{exr:iv09-expand-12}
Can the removable, pole, essential trichotomy be applied at an accumulation point of singularities?
\end{exercise}
\begin{hint}
Check the word isolated.
\end{hint}
\begin{solution}
No. The Laurent annulus around the point may not exist, so the isolated singularity classification does not apply.
\end{solution}

\subsection*{Applications and investigations}

\begin{exercise}[Regularizing a formula]\label{exr:iv09-expand-13}
Why is removing a removable singularity useful in analysis or computation?
\end{exercise}
\begin{hint}
The extended function is genuinely holomorphic at the point.
\end{hint}
\begin{solution}
The extension eliminates artificial division by zero, restores stable local evaluation, and permits direct use of Taylor series and derivative formulas across the point.
\end{solution}

\begin{exercise}[Pole order and blowup]\label{exr:iv09-expand-14}
How does pole order describe the leading local growth?
\end{exercise}
\begin{hint}
Use the factorization \(f=(z-a)^{-m}g\) with \(g(a)\ne0\).
\end{hint}
\begin{solution}
The magnitude behaves like a nonzero constant times \(|z-a|^{-m}\) to leading order, so the integer \(m\) quantifies the algebraic blowup rate.
\end{solution}

\subsection*{Challenge problems}

\begin{exercise}[Reciprocal classification]\label{exr:iv09-expand-15}
If \(f\) has a pole of order \(m\) at \(a\), classify \(1/f\) after extension.
\end{exercise}
\begin{hint}
Invert the local factorization.
\end{hint}
\begin{solution}
Writing \(f=(z-a)^{-m}g\) with \(g(a)\ne0\) gives \(1/f=(z-a)^m/g\), which extends holomorphically with a zero of order \(m\).
\end{solution}

\begin{exercise}[Behavior at infinity]\label{exr:iv09-expand-16}
Explain how to classify the point at infinity for a function by studying \(g(w)=f(1/w)\) at \(w=0\).
\end{exercise}
\begin{hint}
The map \(w=1/z\) turns large \(z\) into small \(w\).
\end{hint}
\begin{solution}
The singularity type of \(g\) at zero defines the type of \(f\) at infinity. For example, a polynomial of degree \(m\) gives a pole of order \(m\) at infinity.
\end{solution}

% END VOL04-EXPANSION IV09-exercises-01
